\chapter{绪论}
\section{研究背景与意义}

人体几何重建旨在从观测数据中恢复具有几何细节的数字人体模型，在医疗康复[1, 2]、混合现实[3, 4]、电影制作[5, 6] 以及远程展示[4, 7] 等领域具有重要价值。传统方法依赖高精度的多相机阵列（如图1­1所示），通过复杂的光学捕捉系统与离线处理流程实现人体数字化。尽管此类方法能生成工业级精度的模型，但对专业设备与场景的严苛要求限制了其在消费级场景中的普及。移动上网设备的普及[8] 推动了数字内容创作与虚拟交互需求的增长，尤其是在远程办公与虚拟社交场景日益普及的背景下，对面向单目图像的人体几何重建技术的需求日益迫切。

这是正文的第一章第一节。根据规范，正文采用小四号宋体，行距为1.5倍\cite{Durham_2017,18,durhamAircraftControlAllocation2017,RN108,RN166,RN172,RN50,RN44,RN40,RN35}。
132
这里测试一下英文和数字的字体：The standard requires Times New Roman for English and numbers, e.g., $E=mc^2$ or 2025.
123

\section{研究内容与贡献}
这是二级标题（小三号黑体）。
123

\section{论文结构安排}
这是三级标题（四号黑体）。
当文字填满一页时，你可以观察页眉在奇数页和偶数页的变化。
